% Cabecera propia de Yava. Se suma a la de los libros (src/agencydomains/_build/header-libro-extra.tex),
% que pone el estilo O'Reilly de cabeceras y pies, la página en blanco con leyenda y los glifos de reserva.
% Yava va en ocho capítulos (las preguntas de primer nivel), cada uno con sus secciones: 1, 1.1, 1.1.1…
\usepackage{float}
\usepackage[figuresright]{rotating}          % sidewaysfigure: las láminas 2 y 5, apaisadas a página completa
\usepackage{needspace}
\usepackage{url}              % las URL de las fuentes se cortan donde haga falta (xurl no está en todas las instalaciones)
\makeatletter\g@addto@macro{\UrlBreaks}{\UrlOrds}\makeatother
\usepackage{xcolor}
\definecolor{vergis}{HTML}{35606F}

% Numeración por capítulo: 1, 1.1, 1.1.1 (la de la clase book); las subsecciones también llevan número
\setcounter{secnumdepth}{2}

% Cabeceras: el capítulo en las pares, la sección en las impares
\fancyhead[LE]{\nouppercase{\small\itshape\leftmark}}
\fancyhead[RO]{\nouppercase{\small\itshape\rightmark}}
\renewcommand{\chaptermark}[1]{\markboth{\thechapter.\ #1}{\thechapter.\ #1}}
\renewcommand{\sectionmark}[1]{\markright{\thesection.\ #1}}
\renewcommand{\subsectionmark}[1]{}

% Títulos: el capítulo abre página, con su número grande en el color de Vergis; la sección, con su número
% en el color de Vergis, pide aire antes de empezar
\usepackage{titlesec}
\titleformat{\chapter}[display]{\normalfont\raggedright}{\fontsize{40}{40}\selectfont\textcolor{vergis}{\thechapter}}{1.2ex}{\huge}
\titleformat{name=\chapter,numberless}{\normalfont\raggedright}{}{0pt}{\huge}
\titlespacing*{\chapter}{0pt}{2ex}{5ex}
\titleformat{\section}{\needspace{8\baselineskip}\normalfont\LARGE\raggedright}{\textcolor{vergis}{\thesection}}{0.6em}{}
\titlespacing*{\section}{0pt}{3.2ex plus 1ex minus .2ex}{1.6ex plus .2ex}
\titleformat{\subsection}{\normalfont\Large\itshape\raggedright}{\thesubsection}{0.6em}{}
\titlespacing*{\subsection}{0pt}{2.6ex plus 1ex minus .2ex}{1.0ex plus .2ex}
\titleformat{\subsubsection}{\normalfont\large\bfseries\raggedright}{\thesubsubsection}{0.6em}{}

% Índice: los capítulos en redonda con aire antes; las secciones con espacio para «8.12»
\makeatletter
\renewcommand*\l@section{\@dottedtocline{1}{1.5em}{2.6em}}
\makeatother

% Tablas: un punto menos, filas con algo de aire
\AtBeginEnvironment{longtable}{\small\renewcommand{\arraystretch}{1.15}}

% Las figuras flotan pero no se adelantan a su sección
\usepackage[section]{placeins}
\renewcommand{\topfraction}{0.9}
\renewcommand{\textfraction}{0.08}
\renewcommand{\floatpagefraction}{0.9}  % una figura de columna alta (la 8, ~0,8 de la caja) va arriba de una página con texto, no sola
\setcounter{totalnumber}{3}

% Párrafos de tabla y texto: menos viudas y huérfanas
\widowpenalty=10000
\clubpenalty=10000

% Holgura para que ninguna línea se salga de la caja (URL largas, palabras largas en negrita)
\setlength{\emergencystretch}{3em}
\Urlmuskip=0mu               % las URL no se estiran («https : / /»); se cortan tras cualquier carácter
\setlength{\headheight}{14pt} % fancyhdr pide 13,6 pt con STIX Two a 11 pt
